%!TEX program = xelatex
\documentclass[11pt, a4paper]{article}

% AA LINGUA E GESTIONE FONT (Specifica per XeLaTeX)
\usepackage{fontspec}
\setmainfont{Latin Modern Roman}
\usepackage[italian]{babel}

% AA LAYOUT DI PAGINA E TIPOGRAFIA (Ottimizzato per Rilegatura)
\usepackage[left=2.4cm, right=1.6cm, top=2cm, bottom=2cm]{geometry}

% AA Spaziatura tra paragrafi controllata (compatta ed evita salti enormi)
\usepackage[skip=3pt plus 1pt minus 1pt]{parskip}
\setlength{\parindent}{0pt}

% AA Microtype per XeLaTeX
\usepackage[protrusion=true]{microtype}
\raggedbottom

% AA GESTIONE DEI COLORI
\usepackage{xcolor}

% BB COLORI E MACRO PER I SIGMA
\definecolor{sigRed}{RGB}{215, 0, 0}
\definecolor{sigBlue}{RGB}{0, 70, 220}
\definecolor{sigGreen}{RGB}{0, 145, 45}

\newcommand{\sigA}{\textcolor{sigRed}{\boldsymbol{\sigma_A}}}
\newcommand{\sigB}{\textcolor{sigBlue}{\boldsymbol{\sigma_B}}}
\newcommand{\sigC}{\textcolor{sigGreen}{\boldsymbol{\sigma_C}}}

% AA MATEMATICA E SIMBOLI
\usepackage{amsmath}
\usepackage{mathtools}
\usepackage{silence}
\usepackage[warnings-off={mathtools-overbracket,mathtools-colon}]{unicode-math}
\setmathfont{Latin Modern Math}

% BB OTTIMIZZAZIONE MATRICI LUNGHE PER LA STAMPA
\setlength{\arraycolsep}{10pt}

% AA Spaziature verticali attorno alle formule matematiche
\setlength{\abovedisplayskip}{3pt plus 1pt minus 1pt}
\setlength{\belowdisplayskip}{3pt plus 1pt minus 1pt}
\setlength{\abovedisplayshortskip}{0pt plus 1pt}
\setlength{\belowdisplayshortskip}{2pt plus 1pt minus 1pt}

% AA GRAFICA ED IMMAGINI
\usepackage{graphicx}
\graphicspath{{./E3_2_Images/}}

% AA PERSONALIZZAZIONE TITOLI
\usepackage{titlesec}
\titleformat{\section}{\normalfont\Large\bfseries\color{blue!80!black}}{\thesection}{1em}{}
\titleformat{name=\section,numberless}{\normalfont\Large\bfseries\color{blue!80!black}}{}{0pt}{}

\titleformat{\subsection}{\normalfont\large\bfseries}{\thesubsection}{1em}{}
\titleformat{name=\subsection,numberless}{\normalfont\large\bfseries}{}{0pt}{}

\titleformat{\subsubsection}{\normalfont\normalsize\bfseries}{\thesubsubsection}{1em}{}
\titleformat{name=\subsubsection,numberless}{\normalfont\normalsize\bfseries}{}{0pt}{}

\titlespacing*{\section}{0pt}{1.2ex plus 0.4ex minus 0.2ex}{0.6ex plus 0.2ex}
\titlespacing*{\subsection}{0pt}{1ex plus 0.3ex minus 0.2ex}{0.5ex plus 0.2ex}
\titlespacing*{\subsubsection}{0pt}{0.8ex plus 0.2ex minus 0.1ex}{0.4ex plus 0.1ex}

\titleformat{\paragraph}[block]{\normalfont\normalsize\bfseries\color{black!80!black}\raggedright}{}{0pt}{}
\titlespacing*{\paragraph}{0pt}{1.2ex plus 0.3ex minus 0.2ex}{0.4ex}

% AA INTEGRAZIONE CODICE MATLAB
\usepackage{matlab}

% BB Riconfigurazione Automatica Titoli Esportati da MATLAB
\renewcommand{\matlabheading}[1]{\section*{#1}}
\renewcommand{\matlabheadingtwo}[1]{\subsection*{#1}}
\renewcommand{\matlabheadingthree}[1]{\subsubsection*{#1}}

% BB Formule E Matrici Simboliche
\renewenvironment{matlabsymbolicoutput}{%
  \par\vspace{2pt}%
  \begingroup
  \centering
}{%
  \par\vspace{2pt}%
  \endgroup
}

% BB Output Testuale Console
\renewenvironment{matlaboutput}{%
  \par\vspace{1pt}%
  \Verbatim[xleftmargin=10pt, formatcom=\color{black!75}, fontsize=\footnotesize]%
}{%
  \endVerbatim
  \vspace{1pt}%
}

% BB Box Codice (Stile Moderno con Barra Laterale Colorata)
\lstset{
  backgroundcolor={},
  aboveskip=0pt,
  belowskip=0pt,
  basicstyle=\ttfamily\footnotesize,
  xleftmargin=0pt
}

\mdfdefinestyle{matlabcode}{%
  topline=false,
  bottomline=false,
  rightline=false,
  leftline=true,
  linewidth=2.5pt,
  linecolor=blue!75!black,
  backgroundcolor=black!3,
  innertopmargin=4pt,
  innerbottommargin=4pt,
  innerleftmargin=9pt,
  innerrightmargin=6pt,
  skipabove=3pt,
  skipbelow=3pt
}

\allowdisplaybreaks

% AA LINK, METADATI ED ELENCHI
\usepackage{hyperref}
\hypersetup{
    colorlinks=true,
    linkcolor=blue!80!black,
    urlcolor=blue!80!black,
    citecolor=blue!80!black,
    pdftitle={Cinematica Diretta Manipolatore SCARA},
    pdfauthor={Francesco Lo Verde},
}
\usepackage{bookmark}

% BB Spaziatura Elenchi (Arieggiata per la Stampa)
\usepackage{enumitem}
\setlist{itemsep=3.5pt plus 1pt minus 1pt, topsep=4pt plus 1pt minus 1pt, partopsep=0pt, parsep=1pt}


\author{Francesco Lo Verde}
\begin{document}

\matlabtitle{Calcolo Cinematica Diretta Con Parametri Di Denavit-Hartenberg}

\begin{par}
\begin{flushleft}
L'obiettivo della cinematica diretta è determinare la posizione e l'orientamento dell'organo terminale (end-effector) a partire dai valori delle variabili di giunto. Si utilizza la convenzione di Denavit-Hartenberg (DH) per modellare la geometria del robot. La convenzione di Denavit-Hartenberg è un metodo standard per assegnare sistemi di riferimento ai bracci di un robot. Per ogni braccio, si definiscono quattro parametri che descrivono la trasformazione geometrica rispetto al braccio precedente.
\end{flushleft}
\end{par}

\begin{par}
\begin{flushleft}
Si inizia definendo le variabili simboliche necessarie per i parametri DH. Questo permette di eseguire i calcoli in forma generale, senza dover specificare subito valori numerici. Successivamente, si assemblano questi parametri in una tabella dove ogni riga corrisponde a un giunto del robot e le colonne rappresentano i parametri $\left(\begin{array}{c}
\theta ,d,\alpha ,a
\end{array}\right)$.
\end{flushleft}
\end{par}

\begin{matlabcode}
clc;
clear;
close all;

% Definizione Variabili Simboliche
syms q1 q2 q3 q4 d1 d2 d3 d4 alpha1 alpha2 alpha3 alpha4 a1 a2 a3 a4 real;

% Tabella Parametri Denavit-Hartenberg (DH)
DH_table = [q1, 0, pi/2, 0;
            q2, 0, 0, d2;
            q3, 0, 0, d3];


% Numero Giunti Manipolatore
n = size(DH_table, 1);
\end{matlabcode}

\begin{par}
\begin{flushleft}
A partire dalla tabella DH, si può calcolare la matrice di trasformazione omogenea $A_i$ per ciascun giunto $i$. Questa matrice, di dimensioni $4\times 4$ e rappresenta la trasformazione (rotazione e traslazione) dal sistema di riferimento del giunto $i$ a quello del giunto $i-1$. Si utilizza la funzione \texttt{Compute\_DH\_Transformations }per generare le matrici simboliche per ogni riga della tabella DH.
\end{flushleft}
\end{par}

\begin{matlabcode}
[T_matrices, R_matrices] = Compute_DH_Transformations(DH_table);

for i = 1:n
    fprintf('Matrice Di Trasformazione Omogenea ^(%d)T_(%d):\n', i-1, i);
    disp(T_matrices(:,:,i));
end
\end{matlabcode}
\begin{matlaboutput}
Matrice Di Trasformazione Omogenea ^(0)T_(1):
\end{matlaboutput}
\begin{matlabsymbolicoutput}
\hskip1em $\displaystyle \left(\begin{array}{cccc}
\cos\left(q_1 \right) & 0 & \sin\left(q_1 \right) & 0\\
\sin\left(q_1 \right) & 0 & -\cos\left(q_1 \right) & 0\\
0 & 1 & 0 & 0\\
0 & 0 & 0 & 1
\end{array}\right)$
\end{matlabsymbolicoutput}
\begin{matlaboutput}
Matrice Di Trasformazione Omogenea ^(1)T_(2):
\end{matlaboutput}
\begin{matlabsymbolicoutput}
\hskip1em $\displaystyle \left(\begin{array}{cccc}
\cos\left(q_2 \right) & -\sin\left(q_2 \right) & 0 & d_2 \,\cos\left(q_2 \right)\\
\sin\left(q_2 \right) & \cos\left(q_2 \right) & 0 & d_2 \,\sin\left(q_2 \right)\\
0 & 0 & 1 & 0\\
0 & 0 & 0 & 1
\end{array}\right)$
\end{matlabsymbolicoutput}
\begin{matlaboutput}
Matrice Di Trasformazione Omogenea ^(2)T_(3):
\end{matlaboutput}
\begin{matlabsymbolicoutput}
\hskip1em $\displaystyle \left(\begin{array}{cccc}
\cos\left(q_3 \right) & -\sin\left(q_3 \right) & 0 & d_3 \,\cos\left(q_3 \right)\\
\sin\left(q_3 \right) & \cos\left(q_3 \right) & 0 & d_3 \,\sin\left(q_3 \right)\\
0 & 0 & 1 & 0\\
0 & 0 & 0 & 1
\end{array}\right)$
\end{matlabsymbolicoutput}
\begin{matlabcode}

T03 = T_matrices(:,:,1) * T_matrices(:,:,2) * T_matrices(:,:,3);
T03 = simplify(T03);

fprintf('Matrice Di Trasformazione Omogenea Totale ^(0)T_(3):\n');
disp(T03);
\end{matlabcode}
\begin{matlaboutput}
Matrice Di Trasformazione Omogenea Totale ^(0)T_(3):
\end{matlaboutput}
\begin{matlabsymbolicoutput}
\hskip1em $\displaystyle \begin{array}{l}
\left(\begin{array}{cccc}
c\left(q_2 +q_3 \right)\,c\left(q_1 \right) & -s\left(q_2 +q_3 \right)\,c\left(q_1 \right) & s\left(q_1 \right) & c\left(q_1 \right)\,\sigma_1 \\
c\left(q_2 +q_3 \right)\,s\left(q_1 \right) & -s\left(q_2 +q_3 \right)\,s\left(q_1 \right) & -c\left(q_1 \right) & s\left(q_1 \right)\,\sigma_1 \\
s\left(q_2 +q_3 \right) & c\left(q_2 +q_3 \right) & 0 & d_3 \,s\left(q_2 +q_3 \right)+d_2 \,s\left(q_2 \right)\\
0 & 0 & 0 & 1
\end{array}\right)\\
\textrm{Dove}\\
\;\;\sigma_1 =d_3 \,c\left(q_2 +q_3 \right)+d_2 \,c\left(q_2 \right)
\end{array}$
\end{matlabsymbolicoutput}

\begin{par}
\begin{flushleft}
La matrice di trasformazione omogenea totale $^1 T_4$ contiene tutte le informazioni sulla posa dell'organo terminale. Da essa possiamo estrarre due elementi chiave:
\end{flushleft}
\end{par}

\begin{itemize}
\setlength{\itemsep}{-1ex}
   \item{\begin{flushleft} La \textbf{matrice di rotazione R}, che definisce l'orientamento dell'organo terminale rispetto alla base. Corrisponde alla sottomatrice 3x3 in alto a sinistra. \end{flushleft}}
   \item{\begin{flushleft} Il \textbf{vettore di posizione p}, che indica la posizione dell'origine del sistema di riferimento dell'organo terminale. Corrisponde alle prime tre righe della quarta colonna. \end{flushleft}}
\end{itemize}

\begin{matlabcode}
% Estrazione Vettore Posizione
p03 = T03(1:3, 4);
p03 = simplify(p03);
fprintf('Vettore Posizione p03:\n');
disp(p03);
\end{matlabcode}
\begin{matlaboutput}
Vettore Posizione p03:
\end{matlaboutput}
\begin{matlabsymbolicoutput}
\hskip1em $\displaystyle \left(\begin{array}{c}
c\left(q_1 \right)\,{\left(d_3 \,c\left(q_2 +q_3 \right)+d_2 \,c\left(q_2 \right)\right)}\\
s\left(q_1 \right)\,{\left(d_3 \,c\left(q_2 +q_3 \right)+d_2 \,c\left(q_2 \right)\right)}\\
d_3 \,s\left(q_2 +q_3 \right)+d_2 \,s\left(q_2 \right)
\end{array}\right)$
\end{matlabsymbolicoutput}
\begin{matlabcode}

% Estrazione Matrice Rotazione
R03 = T03(1:3, 1:3);
R03 = simplify(R03);
fprintf('Matrice Rotazione R03:\n');
disp(R03);
\end{matlabcode}
\begin{matlaboutput}
Matrice Rotazione R03:
\end{matlaboutput}
\begin{matlabsymbolicoutput}
\hskip1em $\displaystyle \left(\begin{array}{ccc}
c\left(q_2 +q_3 \right)\,c\left(q_1 \right) & -s\left(q_2 +q_3 \right)\,c\left(q_1 \right) & s\left(q_1 \right)\\
c\left(q_2 +q_3 \right)\,s\left(q_1 \right) & -s\left(q_2 +q_3 \right)\,s\left(q_1 \right) & -c\left(q_1 \right)\\
s\left(q_2 +q_3 \right) & c\left(q_2 +q_3 \right) & 0
\end{array}\right)$
\end{matlabsymbolicoutput}

\begin{par}
\begin{flushleft}
Per visualizzare una configurazione specifica del robot sono necessari dei valori numerici. Si sostituiscono ora le variabili simboliche (angoli di giunto e parametri geometrici) con valori numerici.
\end{flushleft}
\end{par}
\newpage
\begin{matlabcode}
% Lista Variabili Simboliche
symsList = symvar(T_matrices);
disp('Variabili Simboliche:');
disp(symsList);
\end{matlabcode}

\begin{matlaboutput}
Variabili Simboliche:
\end{matlaboutput}
\begin{matlabsymbolicoutput}
\hskip1em $\displaystyle \left(\begin{array}{ccccc}
d_2  & d_3  & q_1  & q_2  & q_3
\end{array}\right)$
\end{matlabsymbolicoutput}
\begin{matlabcode}

% Valori Numerici Da Sostituire
value = [0.4, 0.3, pi/6, pi/4, -pi/3];

% La sostituzione su un array 3D simbolico può essere problematica.
% È più sicuro sostituire i valori matrice per matrice.
T_Num = zeros(4,4,n);
for i = 1:n
    T_Num(:,:,i) = subs(T_matrices(:,:,i), symsList, value);
end

% Stampa Matrici Numeriche
for i = 1:n
    fprintf('Matrice ^(%d)T_(%d) numerica:\n', i-1, i);
    disp(T_Num(:,:,i));
end
\end{matlabcode}
\begin{matlaboutput}
Matrice ^(0)T_(1) numerica:
    0.8660         0    0.5000         0
    0.5000         0   -0.8660         0
         0    1.0000         0         0
         0         0         0    1.0000
Matrice ^(1)T_(2) numerica:
    0.7071   -0.7071         0    0.2828
    0.7071    0.7071         0    0.2828
         0         0    1.0000         0
         0         0         0    1.0000
Matrice ^(2)T_(3) numerica:
    0.5000    0.8660         0    0.1500
   -0.8660    0.5000         0   -0.2598
         0         0    1.0000         0
         0         0         0    1.0000
\end{matlaboutput}

\matlabheadingtwo{\textbf{Visualizzazione Grafica}}

\begin{matlabcode}
figure;
hold on;
grid on;
xlabel('X');
ylabel('Y');
zlabel('Z');
axis equal;

% Definizione Colori Assi X, Y, Z (standard RGB)
colors = {'r', 'g', 'b'};

% Scala Per Le Frecce Degli Assi
arrow_scale = 0.1;

labels = {'Base', 'Giunto 1', 'Giunto 2', 'End-Effector'};
\end{matlabcode}
\newpage
\begin{matlabcode}
% T_cumulative Tiene Traccia Della Trasformazione Dalla Base Al Frame
% Corrente
T_cumulative = eye(4);
% Memorizza Posizioni Origini Frame
positions = T_cumulative(1:3, 4)';

% Frame Base (RF0)
for j = 1:3
    quiver3(0, 0, 0, T_cumulative(1,j)*arrow_scale, T_cumulative(2,j)*arrow_scale, T_cumulative(3,j)*arrow_scale, ...
        'Color', colors{j}, 'LineWidth', 2, 'MaxHeadSize', 0.4);
end
text(0, 0, 0.01, [' ' labels{1}], 'FontSize', 10);

% Vettori Per Leggenda
h_axes = zeros(1, 3);

% Iterazione Giunti Per Ottenere Posa (da RF1 a RF3)
for i = 1:n
    % Aggiorna Matrice Trasformazione Cumulativa
    T_cumulative = T_cumulative * double(T_Num(:,:,i));

    % Estrazione Posizione E Orientamento
    pos = T_cumulative(1:3, 4);
    R = T_cumulative(1:3, 1:3);

    % Assi Frame Corrente
    for j = 1:3
        h_axes(j) = quiver3(pos(1), pos(2), pos(3), R(1,j)*arrow_scale, R(2,j)*arrow_scale, R(3,j)*arrow_scale, ...
                             'Color', colors{j}, 'LineWidth', 2, 'MaxHeadSize', 0.4);
    end

    % Memorizza Posizione Corrente Per Disegnare Bracci
    positions = [positions; pos'];

    % Etichette
    text(pos(1), pos(2), pos(3) + 0.01, [' ' labels{i+1}], 'FontSize', 10);
end

% Disegna I Bracci collegando le origini dei frame memorizzate
h_line = plot3(positions(:,1), positions(:,2), positions(:,3), 'k-o', 'LineWidth', 1.5, 'MarkerSize', 5, 'MarkerFaceColor', 'k');

legend([h_axes, h_line], {'Asse X', 'Asse Y', 'Asse Z', 'Bracci'}, 'Location', 'Best');
title('Configurazione Manipolatore Antropomorfo');
view(150, 20);

axis tight;
hold off;
\end{matlabcode}
\begin{center}
\includegraphics[width=1\linewidth]{figure_0.png}
\end{center}

\end{document}
